\documentclass[
  12pt,
  oneside,
  a4paper,
  brazil
]{abntex2}

\usepackage[utf8]{inputenc}
\usepackage[T1]{fontenc}
\usepackage{lmodern}
\usepackage{babel}
\usepackage{microtype}
\usepackage{graphicx}
\usepackage{booktabs}
\usepackage{longtable}
\usepackage{tabularx}
\usepackage{array}
\usepackage{xcolor}
\usepackage{hyperref}
\usepackage{enumitem}
\usepackage{indentfirst}
\usepackage{geometry}

\geometry{left=3cm,right=2cm,top=3cm,bottom=2cm}
\setlength{\parindent}{1.25cm}
\setlength{\parskip}{0.2cm}
\renewcommand{\arraystretch}{1.25}
\newcolumntype{Y}{>{\raggedright\arraybackslash}X}

\hypersetup{
  colorlinks=true,
  linkcolor=black,
  citecolor=black,
  urlcolor=blue,
  pdftitle={AgroSmart+ - Documentação Técnica e Acadêmica},
  pdfauthor={Equipe AgroSmart+}
}

\titulo{AgroSmart+: Plataforma Web para Apoio ao Agronegócio}
\autor{Equipe AgroSmart+}
\local{Brasil}
\data{2026}
\instituicao{%
  Instituição de Ensino\\
  Curso de Desenvolvimento de Sistemas ou área correlata
}
\tipotrabalho{Documentação técnica e acadêmica}
\preambulo{Documentação técnica e acadêmica apresentada como material de avaliação do projeto AgroSmart+, contemplando análise de requisitos, arquitetura, implementação, segurança, responsividade, banco de dados, integração com APIs externas e planejamento de evolução.}

\begin{document}
\selectlanguage{brazil}
\frenchspacing

\imprimircapa
\imprimirfolhaderosto

\begin{resumo}
O AgroSmart+ é uma plataforma web acadêmica voltada ao agronegócio, desenvolvida com frontend em React e Vite, backend em Node.js com Express e banco de dados PostgreSQL. A solução reúne funcionalidades de cadastro e autenticação de usuários, consulta de cotações agrícolas, previsão do tempo, FAQ, consulta técnica baseada na API RespondeAgro da Embrapa, recursos de acessibilidade e painel administrativo. Esta documentação apresenta a análise técnica do projeto, sua arquitetura, requisitos funcionais e não funcionais, modelagem de dados, fluxos de autenticação, funcionamento do reCAPTCHA, integrações externas, segurança, responsividade, limitações atuais e possibilidades de evolução. O objetivo é fornecer uma visão completa do sistema em nível adequado para banca avaliadora, mantendo fidelidade ao estado real do código-fonte.

\textbf{Palavras-chave}: Agronegócio. Desenvolvimento Web. React. Node.js. PostgreSQL. Sistemas acadêmicos.
\end{resumo}

\pdfbookmark[0]{\contentsname}{toc}
\tableofcontents*
\cleardoublepage

\textual

\chapter{Introdução}

O agronegócio brasileiro depende cada vez mais de informações digitais para tomada de decisão. Dados sobre clima, preço de commodities, localização, histórico de cotações, suporte técnico e gestão de perfis costumam estar distribuídos em serviços distintos, o que dificulta o acesso rápido e organizado por produtores, empresas e cooperativas.

O AgroSmart+ surge nesse contexto como uma plataforma web acadêmica que centraliza funcionalidades úteis ao setor agrícola. O sistema combina uma interface moderna e responsiva com uma API própria, persistência em PostgreSQL e integração com serviços externos. A proposta do projeto é demonstrar domínio de desenvolvimento full stack, organização de arquitetura, segurança básica, documentação técnica e adequação visual ao domínio do agronegócio.

\section{Identificação do Projeto}

\begin{table}[h]
\centering
\caption{Identificação do projeto}
\begin{tabularx}{\textwidth}{@{}lY@{}}
\toprule
\textbf{Item} & \textbf{Descrição} \\
\midrule
Nome & AgroSmart+ \\
Área & Desenvolvimento Web aplicado ao agronegócio \\
Tipo & Projeto acadêmico full stack \\
Frontend & React, Vite e Tailwind CSS \\
Backend & Node.js, Express e PostgreSQL \\
Banco de dados & PostgreSQL local ou Supabase PostgreSQL \\
Deploy previsto & Vercel, Render e Supabase \\
Público & Agricultores, empresários agrícolas, cooperativas e administradores \\
\bottomrule
\end{tabularx}
\end{table}

\section{Justificativa}

A justificativa do projeto está relacionada à necessidade de reunir, em uma única solução, informações agrícolas relevantes que normalmente exigem navegação por diferentes plataformas. O AgroSmart+ contribui para o aprendizado acadêmico ao integrar conceitos de frontend, backend, banco de dados, autenticação, consumo de APIs externas, segurança, responsividade e documentação formal.

\section{Objetivo Geral}

Desenvolver uma plataforma web para apoio ao agronegócio, permitindo cadastro de usuários, autenticação segura, consulta de cotações, clima, perguntas frequentes, informações técnicas e administração de dados do sistema.

\section{Objetivos Específicos}

\begin{itemize}
  \item Implementar cadastro de usuários por tipo de perfil.
  \item Autenticar usuários e administradores com JWT.
  \item Proteger senhas de usuários com bcrypt.
  \item Consultar e exibir cotações agrícolas em dashboard responsivo.
  \item Integrar consulta climática via OpenWeatherMap.
  \item Integrar consulta técnica via API RespondeAgro da Embrapa.
  \item Disponibilizar painel administrativo para usuários, FAQ e tabelas.
  \item Garantir compatibilidade com deploy em Vercel, Render e Supabase.
  \item Documentar variáveis de ambiente, arquitetura, requisitos e segurança.
\end{itemize}

\section{Problema Identificado}

Usuários do agronegócio precisam acessar dados de mercado, clima e suporte técnico de forma rápida, mas frequentemente encontram informações dispersas. Essa fragmentação aumenta o tempo de consulta e dificulta decisões operacionais. O projeto aborda esse problema por meio de uma interface única, com autenticação e módulos especializados.

\section{Público-alvo}

O público-alvo do AgroSmart+ inclui produtores rurais, empresários do setor agrícola, cooperativas e administradores da plataforma. No contexto acadêmico, o projeto também se destina a avaliadores, professores e estudantes interessados em arquitetura full stack.

\section{Escopo do Sistema}

O escopo atual contempla frontend web, API backend, banco PostgreSQL, autenticação JWT, cadastro com verificação de e-mail, painel de cotações, clima, FAQ, RespondeAgro, painel administrativo e página 404 personalizada. Não fazem parte do escopo atual: aplicativo mobile nativo, pagamento, BI avançado, auditoria completa, multi-tenant corporativo e automação agrícola em tempo real.

\chapter{Tecnologias Utilizadas}

\begin{longtable}{@{}p{0.22\textwidth}p{0.70\textwidth}@{}}
\caption{Tecnologias do projeto}\\
\toprule
\textbf{Categoria} & \textbf{Tecnologias e uso} \\
\midrule
\endfirsthead
\toprule
\textbf{Categoria} & \textbf{Tecnologias e uso} \\
\midrule
\endhead
Frontend & React 19, Vite, React Router, Tailwind CSS, DaisyUI, Axios, React Hook Form, Lucide React, React Icons, Recharts e SweetAlert2. \\
Backend & Node.js, Express 5, CORS, dotenv, PostgreSQL com \texttt{pg}, bcrypt, jsonwebtoken, Nodemailer, Axios, Swagger UI e Puppeteer. \\
Banco & PostgreSQL, com compatibilidade prevista para Supabase PostgreSQL. \\
Autenticação & JWT para sessões e bcrypt para hash de senhas de usuários. \\
CAPTCHA & Google reCAPTCHA, com site key no frontend e secret key no backend. \\
APIs externas & IBGE Localidades, OpenWeatherMap, Embrapa RespondeAgro, SendGrid SMTP/Nodemailer e páginas de cooperativas para cotações. \\
Deploy & Vercel para frontend, Render para backend e Supabase para banco. \\
\bottomrule
\end{longtable}

\chapter{Arquitetura do Sistema}

A arquitetura do AgroSmart+ é separada em três camadas principais: frontend, backend e banco de dados. O frontend é responsável pela experiência do usuário, navegação, formulários e visualização dos dados. O backend concentra regras de negócio, autenticação, acesso ao banco, validações, envio de e-mails, integração com APIs externas e scrapers. O banco PostgreSQL persiste usuários, perfis, cotações, FAQ, administradores e informações auxiliares.

\section{Visão Geral da Arquitetura}

\begin{verbatim}
Usuário
  |
  v
Frontend React/Vite
  |
  v
Backend Node.js/Express
  |
  v
PostgreSQL/Supabase
  |
  +--> IBGE Localidades
  +--> OpenWeatherMap
  +--> Embrapa RespondeAgro
  +--> SendGrid SMTP/Nodemailer
  +--> Scrapers de cotações
\end{verbatim}

% Inserir aqui um diagrama real de arquitetura, se disponível.
% Exemplo:
% \begin{figure}[h]
% \centering
% \includegraphics[width=0.9\textwidth]{figuras/arquitetura.png}
% \caption{Arquitetura geral do AgroSmart+}
% \end{figure}

\section{Estrutura de Pastas}

\begin{verbatim}
AgroSmart/
  backend/
    routes/
    services/
    utils/
    server.js
    supabase-schema.sql
  frontend/
    src/
      assets/
      componentes/
      contexts/
      hooks/
      pages/
      services/
  docs/
    agrosmart-documentacao.tex
  render.yaml
  DEPLOY.md
  package.json
  readme.md
\end{verbatim}

\chapter{Requisitos}

\section{Requisitos Funcionais}

\begin{longtable}{@{}p{0.12\textwidth}p{0.72\textwidth}p{0.12\textwidth}@{}}
\caption{Requisitos funcionais}\\
\toprule
\textbf{Código} & \textbf{Descrição} & \textbf{Status} \\
\midrule
\endfirsthead
\toprule
\textbf{Código} & \textbf{Descrição} & \textbf{Status} \\
\midrule
\endhead
RF01 & Permitir cadastro de usuários com perfis de agricultor, empresário e cooperativa. & Implementado \\
RF02 & Validar CPF e CNPJ no cadastro. & Implementado \\
RF03 & Enviar código de verificação por e-mail durante o cadastro. & Implementado \\
RF04 & Autenticar usuários com e-mail e senha. & Implementado \\
RF05 & Gerar token JWT após login válido. & Implementado \\
RF06 & Proteger rotas internas no frontend. & Implementado \\
RF07 & Proteger rotas sensíveis no backend com JWT. & Implementado \\
RF08 & Consultar cotações agrícolas e histórico. & Implementado \\
RF09 & Consultar previsão do tempo. & Implementado \\
RF10 & Permitir envio de mensagens pelo FAQ. & Implementado \\
RF11 & Consultar informações técnicas pela Embrapa. & Implementado \\
RF12 & Listar e alterar status de usuários no painel administrativo. & Implementado \\
RF13 & Consultar tabelas do banco pelo painel administrativo. & Implementado \\
RF14 & Criar novos administradores pelo painel. & Previsto \\
RF15 & Exibir estatísticas administrativas avançadas. & Em evolução \\
\bottomrule
\end{longtable}

\section{Requisitos Não Funcionais}

\begin{longtable}{@{}p{0.12\textwidth}p{0.72\textwidth}p{0.12\textwidth}@{}}
\caption{Requisitos não funcionais}\\
\toprule
\textbf{Código} & \textbf{Descrição} & \textbf{Status} \\
\midrule
\endfirsthead
\toprule
\textbf{Código} & \textbf{Descrição} & \textbf{Status} \\
\midrule
\endhead
RNF01 & Interface responsiva para desktop, tablet e mobile. & Implementado \\
RNF02 & Separação entre frontend, backend e banco. & Implementado \\
RNF03 & Uso de variáveis de ambiente para credenciais. & Implementado \\
RNF04 & Senhas de usuários protegidas com bcrypt. & Implementado \\
RNF05 & Autenticação baseada em JWT. & Implementado \\
RNF06 & Compatibilidade com Supabase, Render e Vercel. & Implementado \\
RNF07 & Tratamento de erros sem expor stack trace em produção. & Implementado \\
RNF08 & Documentação técnica e acadêmica. & Implementado \\
RNF09 & Testes automatizados abrangentes. & Previsto \\
\bottomrule
\end{longtable}

\section{Regras de Negócio}

\begin{itemize}
  \item Cada usuário deve possuir um tipo: agricultor, empresário ou cooperativa.
  \item O e-mail do usuário deve ser único.
  \item Agricultores e empresários podem informar CPF.
  \item Empresários e cooperativas podem informar CNPJ.
  \item Senhas de usuários comuns devem ser armazenadas com hash.
  \item Usuários autenticados recebem JWT com validade limitada.
  \item Rotas administrativas exigem autenticação.
  \item O reCAPTCHA somente é obrigatório quando \texttt{CAPTCHA\_ENABLED=true}.
\end{itemize}

\chapter{Casos de Uso e Histórias de Usuário}

\section{Casos de Uso}

\begin{longtable}{@{}p{0.16\textwidth}p{0.24\textwidth}p{0.52\textwidth}@{}}
\caption{Casos de uso principais}\\
\toprule
\textbf{Código} & \textbf{Ator} & \textbf{Caso de uso} \\
\midrule
\endfirsthead
\toprule
\textbf{Código} & \textbf{Ator} & \textbf{Caso de uso} \\
\midrule
\endhead
UC01 & Visitante & Criar conta informando dados pessoais, localização e perfil. \\
UC02 & Usuário & Fazer login e acessar dashboard protegido. \\
UC03 & Usuário & Consultar cotações de grãos e histórico. \\
UC04 & Usuário & Consultar clima por cidade ou localização. \\
UC05 & Usuário & Enviar dúvida pelo FAQ. \\
UC06 & Usuário & Consultar base técnica da Embrapa. \\
UC07 & Administrador & Listar usuários e alterar status. \\
UC08 & Administrador & Consultar tabelas do banco. \\
\bottomrule
\end{longtable}

\section{Histórias de Usuário}

\begin{itemize}
  \item Como agricultor, quero consultar cotações para acompanhar oportunidades de venda.
  \item Como empresário agrícola, quero visualizar preços e clima para apoiar decisões comerciais.
  \item Como cooperativa, quero acessar dados do sistema para acompanhar usuários e demandas.
  \item Como administrador, quero gerenciar status de usuários para manter a integridade da plataforma.
  \item Como visitante, quero criar uma conta de forma guiada e segura.
\end{itemize}

\chapter{Modelagem do Banco de Dados}

O banco de dados utiliza PostgreSQL. O arquivo \texttt{backend/supabase-schema.sql} contém o schema principal, com tabelas para usuários, perfis, grãos, relações usuário-grão, cotações, FAQ e administradores.

\section{Modelo Entidade-Relacionamento}

As principais entidades são: usuário, agricultor, empresário, cooperativa, grão, usuário-grão, histórico de cotação, FAQ, administrador, cache de cotações e histórico de cotações coletadas por scrapers.

% Inserir aqui uma imagem real do MER/DER, se disponível.
% Exemplo:
% \begin{figure}[h]
% \centering
% \includegraphics[width=0.95\textwidth]{figuras/mer-agrosmart.png}
% \caption{Modelo Entidade-Relacionamento do AgroSmart+}
% \end{figure}

\section{Dicionário de Dados}

\begin{longtable}{@{}p{0.22\textwidth}p{0.22\textwidth}p{0.42\textwidth}@{}}
\caption{Dicionário de dados resumido}\\
\toprule
\textbf{Tabela} & \textbf{Campo} & \textbf{Descrição} \\
\midrule
\endfirsthead
\toprule
\textbf{Tabela} & \textbf{Campo} & \textbf{Descrição} \\
\midrule
\endhead
tb\_usuario & id & Identificador primário do usuário. \\
tb\_usuario & nome\_completo & Nome completo informado no cadastro. \\
tb\_usuario & email & E-mail único usado para login. \\
tb\_usuario & senha & Hash da senha do usuário. \\
tb\_usuario & tipo\_usuario & Perfil: agricultor, empresario ou cooperativa. \\
tb\_usuario & status & Estado ativo ou inativo. \\
tb\_agricultor & usuario\_id & Referência ao usuário. \\
tb\_agricultor & cpf & CPF do agricultor. \\
tb\_agricultor & nome\_propriedade & Nome da propriedade rural. \\
tb\_empresario & cnpj & CNPJ da empresa, quando informado. \\
tb\_cooperativa & nome\_cooperativa & Nome da cooperativa. \\
tb\_grao & nome & Nome do grão. \\
tb\_usuario\_grao & tipo\_relacao & Relação cultiva ou interesse. \\
tb\_faq & mensagem & Mensagem enviada pelo formulário de suporte. \\
tb\_admin & email & E-mail do administrador. \\
tb\_cotacoes\_cache & dados & JSON com cotações cacheadas por provedor. \\
tb\_cotacoes\_historico & preco & Preço coletado para histórico de mercado. \\
\bottomrule
\end{longtable}

\chapter{Descrição das Principais Telas}

\section{Home}

A página inicial apresenta a identidade visual do AgroSmart+, com paleta verde, elementos agrícolas, chamadas para login/cadastro, apresentação de soluções e parceiros. Ela serve como referência visual para páginas públicas, incluindo a página 404.

\section{Cadastro}

A tela de cadastro é organizada em etapas. O usuário informa dados pessoais, estado, cidade, tipo de perfil e campos específicos. O sistema pode exigir reCAPTCHA quando habilitado e envia código de verificação por e-mail antes de finalizar o registro.

\section{Login}

A tela de login autentica usuários e administradores. O frontend envia e-mail, senha e token do reCAPTCHA quando configurado. O backend valida credenciais, gera JWT e retorna os dados do usuário sem senha.

\section{Dashboard de Cotações}

O dashboard exibe cooperativas, grãos, preços atuais, filtros e histórico. A tela foi estruturada para funcionar em desktop e telas menores, com tabelas em contêineres roláveis quando necessário.

\section{Clima}

A página de clima consulta a API OpenWeatherMap. A chave é configurada por \texttt{VITE\_OPENWEATHER\_API\_KEY}. Caso a chave não esteja presente, a interface informa erro amigável sem quebrar a aplicação.

\section{FAQ e RespondeAgro}

O FAQ permite envio de mensagens ao banco. A tela RespondeAgro consulta a API da Embrapa por meio do backend, evitando que credenciais sensíveis fiquem expostas no frontend.

\section{Painel Administrativo}

O painel administrativo possui dashboard, gerenciamento de usuários, consulta de tabelas, mensagens do FAQ e páginas previstas para estatísticas, logs e criação de administradores. Algumas funcionalidades administrativas ainda estão em evolução.

\section{Página 404}

A página 404 segue a identidade visual da Home, reutilizando navbar, footer, fundo agrícola e uma ilustração de colheitadeira com falha. Ela é carregada por rota fallback do React Router.

% Inserir aqui prints reais das telas.
% Exemplo:
% \includegraphics[width=\textwidth]{figuras/home.png}

\chapter{Fluxos do Sistema}

\section{Fluxo de Cadastro e Login}

No cadastro, o usuário preenche os dados em etapas, seleciona estado/cidade, escolhe perfil e informa campos específicos. Se o CAPTCHA estiver habilitado, o frontend obtém um token do Google reCAPTCHA. O backend valida esse token antes de aceitar o registro. Após a verificação de e-mail, o usuário é persistido em transação no PostgreSQL.

No login, o usuário envia e-mail, senha e, quando aplicável, token de reCAPTCHA. O backend consulta usuário ou administrador, valida senha, gera JWT e retorna o token para o frontend.

\section{Fluxo de Autenticação JWT}

\begin{enumerate}
  \item O usuário realiza login.
  \item O backend valida as credenciais.
  \item O backend gera JWT com identificador, e-mail e tipo de usuário.
  \item O frontend armazena token e dados do usuário no \texttt{localStorage}.
  \item O cliente Axios envia \texttt{Authorization: Bearer <token>}.
  \item O middleware do backend valida o token nas rotas protegidas.
\end{enumerate}

\section{Funcionamento do reCAPTCHA}

O reCAPTCHA é controlado por variáveis de ambiente. A site key é pública e usada no frontend. A secret key é privada e usada apenas pelo backend. Quando \texttt{CAPTCHA\_ENABLED=false}, o backend não exige token, permitindo desenvolvimento local sem travamento.

\chapter{Integrações Externas}

\begin{longtable}{@{}p{0.25\textwidth}p{0.62\textwidth}p{0.10\textwidth}@{}}
\caption{Integrações externas}\\
\toprule
\textbf{Integração} & \textbf{Uso} & \textbf{Status} \\
\midrule
\endfirsthead
\toprule
\textbf{Integração} & \textbf{Uso} & \textbf{Status} \\
\midrule
\endhead
IBGE Localidades & Estados e municípios no cadastro. & Ativa \\
OpenWeatherMap & Clima atual e previsão. & Ativa \\
Embrapa RespondeAgro & Consulta técnica pelo backend. & Ativa \\
SendGrid SMTP/Nodemailer & Envio de códigos de verificação e senha pela porta 2525 no Render Free. & Ativa \\
Google reCAPTCHA & Proteção de login e cadastro quando habilitada. & Ativa \\
Scrapers de cotações & Coleta de preços por provedores como COAMO, LAR e Granos. & Em evolução \\
\bottomrule
\end{longtable}

\chapter{Segurança do Sistema}

O AgroSmart+ aplica medidas básicas de segurança compatíveis com o estágio acadêmico do projeto:

\begin{itemize}
  \item Senhas de usuários comuns são armazenadas com bcrypt.
  \item JWT usa segredo vindo de \texttt{JWT\_SECRET}.
  \item Rotas protegidas usam middleware de autenticação.
  \item Arquivos \texttt{.env} são ignorados pelo Git.
  \item reCAPTCHA é validado no backend quando habilitado.
  \item Chaves sensíveis de Embrapa, e-mail e banco ficam no backend.
  \item Erros internos não expõem stack trace em produção.
  \item CORS restringe origens por \texttt{FRONTEND\_URL}.
\end{itemize}

Como limitação, administradores legados podem existir com senha demonstrativa em scripts de exemplo. Para produção real, recomenda-se migrar administradores para bcrypt, aplicar autorização granular de papel administrativo, rate limiting, auditoria, logs estruturados e rotação de credenciais.

\chapter{Responsividade e Acessibilidade}

O frontend utiliza Tailwind CSS e componentes responsivos. Foram observados e corrigidos pontos de risco visual em autenticação, navbar, tabelas e painéis administrativos. A navbar possui menu mobile, formulários evitam altura fixa em dispositivos pequenos e tabelas usam rolagem horizontal controlada quando não há espaço suficiente.

O projeto também possui recursos de acessibilidade, como menu de acessibilidade, modo escuro, alto contraste, ajuste de fonte e suporte a comandos de voz. Esses recursos devem ser testados continuamente para garantir contraste, foco visível e navegação por teclado.

\chapter{Testes e Validação}

O projeto possui scripts de execução, build e lint em \texttt{package.json}. A validação recomendada inclui:

\begin{itemize}
  \item \texttt{npm run install:all}
  \item \texttt{npm run dev}
  \item \texttt{npm run build}
  \item \texttt{npm run lint}
  \item \texttt{GET /health}
  \item \texttt{GET /test/db-status}
  \item Teste manual de cadastro com \texttt{CAPTCHA\_ENABLED=false}
  \item Teste manual de cadastro com \texttt{CAPTCHA\_ENABLED=true}
  \item Teste manual de login, dashboard, clima, FAQ, painel admin e página 404
\end{itemize}

Até o momento não foi identificada suíte automatizada completa de testes. A criação desses testes é recomendada como evolução.

\chapter{Deploy}

O deploy previsto usa Supabase para banco, Render para backend e Vercel para frontend.

\section{Supabase}

Executar \texttt{backend/supabase-schema.sql} no SQL Editor. Configurar a connection string no Render como \texttt{DATABASE\_URL}. Em cenários com SSL, usar connection string com SSL ou \texttt{PGSSL=require}.

\section{Render}

O arquivo \texttt{render.yaml} define serviço web Node com diretório raiz \texttt{backend}, build \texttt{npm install}, start \texttt{npm start} e health check em \texttt{/health}. Variáveis sensíveis devem ser configuradas no painel do Render.

\section{Vercel}

O frontend deve ser configurado com diretório raiz \texttt{frontend}, build \texttt{npm run build} e output \texttt{dist}. O arquivo \texttt{frontend/vercel.json} garante rewrite para rotas do React Router.

\chapter{Cronograma}

\begin{table}[h]
\centering
\caption{Cronograma acadêmico sugerido}
\begin{tabularx}{\textwidth}{@{}lYl@{}}
\toprule
\textbf{Etapa} & \textbf{Atividade} & \textbf{Status} \\
\midrule
1 & Levantamento de requisitos e definição do escopo. & Concluído \\
2 & Desenvolvimento do frontend e identidade visual. & Concluído \\
3 & Desenvolvimento do backend e banco de dados. & Concluído \\
4 & Integrações externas e cotações. & Em evolução \\
5 & Responsividade e ajustes de segurança. & Concluído \\
6 & Documentação técnica e acadêmica. & Concluído \\
7 & Testes automatizados e hardening para produção. & Previsto \\
\bottomrule
\end{tabularx}
\end{table}

\chapter{Resultados Obtidos}

O AgroSmart+ apresenta uma aplicação web full stack com identidade visual coerente ao agronegócio, autenticação, cadastro por perfil, dashboard de cotações, clima, FAQ, consulta técnica e administração. A estrutura permite implantação em serviços modernos de nuvem e está documentada para continuidade acadêmica.

\chapter{Limitações Atuais}

\begin{itemize}
  \item Nem todas as rotas estão documentadas no Swagger.
  \item Algumas páginas administrativas são demonstrativas ou estão em evolução.
  \item A edição completa de perfil ainda requer persistência ampliada.
  \item Não há suíte automatizada de testes abrangente.
  \item Autorização administrativa pode ser refinada para checagem de papel em todas as ações sensíveis.
  \item Scrapers dependem da estrutura de páginas externas, podendo exigir manutenção.
\end{itemize}

\chapter{Trabalhos Futuros}

\begin{itemize}
  \item Implementar CRUD completo de FAQ administrativo.
  \item Criar cadastro seguro de administradores com bcrypt.
  \item Adicionar testes automatizados de frontend, backend e banco.
  \item Expandir Swagger/OpenAPI.
  \item Implementar logs estruturados e auditoria.
  \item Adicionar rate limiting e proteção contra abuso.
  \item Criar diagramas reais de arquitetura e MER.
  \item Aprimorar analytics e estatísticas administrativas.
\end{itemize}

\chapter{Conclusão}

O AgroSmart+ atende à proposta acadêmica de uma plataforma full stack aplicada ao agronegócio, demonstrando integração entre frontend moderno, backend Express, banco PostgreSQL, autenticação JWT, validações, integrações externas e preocupação com responsividade e documentação. Apesar de existirem limitações compatíveis com o estágio do projeto, a base atual é consistente e permite evolução para cenários mais robustos de uso, avaliação e demonstração.

\postextual

\chapter*{Referências}
\addcontentsline{toc}{chapter}{Referências}

ASSOCIAÇÃO BRASILEIRA DE NORMAS TÉCNICAS. \textbf{NBR 14724: Informação e documentação: trabalhos acadêmicos: apresentação}. Rio de Janeiro: ABNT, 2011.

ASSOCIAÇÃO BRASILEIRA DE NORMAS TÉCNICAS. \textbf{NBR 6023: Informação e documentação: referências: elaboração}. Rio de Janeiro: ABNT, 2018.

NODE.JS. \textbf{Node.js documentation}. Disponível em: \url{https://nodejs.org/}. Acesso em: 27 maio 2026.

REACT. \textbf{React documentation}. Disponível em: \url{https://react.dev/}. Acesso em: 27 maio 2026.

EXPRESS. \textbf{Express documentation}. Disponível em: \url{https://expressjs.com/}. Acesso em: 27 maio 2026.

POSTGRESQL. \textbf{PostgreSQL documentation}. Disponível em: \url{https://www.postgresql.org/docs/}. Acesso em: 27 maio 2026.

SUPABASE. \textbf{Supabase documentation}. Disponível em: \url{https://supabase.com/docs}. Acesso em: 27 maio 2026.

VERCEL. \textbf{Vercel documentation}. Disponível em: \url{https://vercel.com/docs}. Acesso em: 27 maio 2026.

RENDER. \textbf{Render documentation}. Disponível em: \url{https://render.com/docs}. Acesso em: 27 maio 2026.

GOOGLE. \textbf{reCAPTCHA documentation}. Disponível em: \url{https://developers.google.com/recaptcha}. Acesso em: 27 maio 2026.

EMBRAPA. \textbf{APIs e serviços digitais da Embrapa}. Disponível em: \url{https://www.embrapa.br/}. Acesso em: 27 maio 2026.

\end{document}
